\subsection{Phase-space measure of the invariant family}
\label{sec:measure-definition}

The positive-measure family established above motivates a domain-specific
measure of quasiperiodic initial states. A quantitative fraction requires
a specified volume measure and initial-state domain. Following the measure formulation of
\citet{biasco2018measure,chierchia2021measure}, let $U\subset\R^4$ be a
bounded domain of positive Liouville volume for the identified
two-degree-of-freedom Hamiltonian $h$, and
let $\mathcal K_h$ denote the union of its analytic Lagrangian invariant
tori carrying Diophantine quasiperiodic motion. Define
\begin{equation}
 f_{\rm KAM}(U)=\frac{\mu_L(U\cap\mathcal K_h)}{\mu_L(U)},
 \qquad \mu_L=\mathrm dX_1\,\mathrm dX_2\,\mathrm dY_1\,\mathrm dY_2.
 \label{eq:measure-fraction-definition}
\end{equation}
An explicitly certified subfamily provides a lower bound for this
fraction. The complement of that subfamily can contain other regular
motion; it is not thereby a chaotic set. The sampling confidence in a
learning bound such as that of \citet{chen2022} is a probability over
training data, whereas \eqref{eq:measure-fraction-definition} measures
initial states of one fixed fitted Hamiltonian.

For a symplectic chart $\Xi$ and an action box $D$, set
$U=\Xi(\mathbb T^2_{2\pi}\times D)$, where
$\mathbb T^2_{2\pi}=(\R/2\pi\mathbb Z)^2$. Then
$\mu_L(U)=(2\pi)^2\operatorname{area}(D)$.
The period-one angles used by the earlier torus certificate are changed
to radians together with the inverse action rescaling, preserving the
symplectic form. The constant volume factor in the SI canonical chart
also cancels from the relative fraction. The radius and location of $D$
must accompany every reported percentage.

Where the assumptions of Theorem~\ref{thm:full-fe-kam-lift} hold, the
invariant-family lift to the $64$-mode closed loop carries this
low-dimensional parameter measure forward. It does not turn it into a
positive fraction of the full $128$-dimensional Lebesgue volume.
Accordingly, the measure estimate in Appendix~\ref{app:measure-proof}
and the numerical sampling experiment in Section~\ref{sec:measure-numerical}
refer to the conservative low-mode Hamiltonian realized by the magnetic
feedback, with their respective domains specified separately.


The explicit constants of \citet{chierchia2021measure} yield a computable
local measure bound after construction of a symplectic tube and estimation
of the full analytic remainder. Appendix~\ref{app:measure-proof} gives that
construction for the integral-COMSOL Hamiltonian. Its sufficient estimates
certify an extremely thin transverse domain; the central motion remains
at millimetre amplitude, but the resulting width does not quantify an
experimentally resolvable initial-state tolerance. To examine a broader
range of amplitudes, Section~\ref{sec:measure-numerical} samples a prescribed
action domain and evaluates finite-time quasiperiodicity with independent
integration and averaging checks.
